%
% File: chap05.tex
% Author: Solal Rapaport
% Description: Secret removal in history-altered repositories. Grounded in
% contributions/secrets_removal: implemented detection pipeline, current early
% results, validation protocol, and unfinished perspectives.
%
\chapter{Secret Removal in Version Control Archives}
\label{chap:secrets}

\section{Introduction}
\label{sec:sec-intro}

Chapter~\ref{chap:histories} ended with a specific gap left by the study of altered Git histories.
It showed that secret-shaped files can be removed through history rewriting and still remain recoverable from Software Heritage, but archival recovery alone cannot answer the security question that matters after a leak: whether the underlying credential is still valid and usable.
This chapter starts from that distinction.
It treats secret removal not as another instance of repository hygiene, but as a security response whose effectiveness depends on an object outside Git: the credential itself.

The context is familiar to anyone who has accidentally pushed a password, API token or a private key to a public repository.
Once the mistake is noticed, rewriting the repository history is an available and visible action: the developer can force-push, filter the offending file or lines out of past commits, and make the current repository state look clean again.
That action can be useful as cleanup.
It reduces the chance that a future clone, search, or casual inspection of the repository will rediscover the value.
It may also be required by organisational policy or by the desire not to keep known-sensitive material in the advertised project history.
But it is not the security remediation.
The only action that invalidates the leaked capability is rotation: revoking the exposed token, replacing the key, resetting the password, or otherwise making the old value unusable.
\NOTEside{interpretation}
History rewriting acts on the repository record; rotation acts on the authority that the leaked value grants.

The problem is that these two actions can be confused.
Chapter~\ref{chap:histories} already showed that secret-shaped removals are not isolated curiosities: enough private-key, credential-file, and authentication-token removals were visible in altered histories to make secret suppression a first-order case study rather than an anecdote.
That evidence raises a behavioural and operational question.
If developers invest effort in rewriting history after a leak, do some of them treat that rewrite as the solution rather than as cleanup after rotation?
The archive cannot answer that question directly, but it can show whether the secrets that developers tried to remove remain recoverable.
If they do, and if even a fraction of them are still valid, the consequence is not merely an untidy Git history.
Valid secrets can open access to source repositories, deployment systems, cloud accounts, payment providers, messaging platforms, databases, or continuous-integration environments, depending on the credential family.

Some platforms reduce this risk by scanning pushed content for known secret formats and by warning users, notifying providers, or invalidating specific token families.
Those mechanisms are valuable because they act on the credential or alert the actor who can revoke it.
They are also necessarily partial.
They depend on the platform that receives the push, the detector coverage, the provider integration, and the type of credential involved.
Private keys, database credentials, self-hosted services, certificates, copied blobs, archives, forks, and third-party mirrors can fall outside the cleanest version of this model.
The result is a gap between repository-level deletion and credential-level remediation: the former can be observed in Git history, while the latter must be verified against the authority that accepts the credential.

This chapter addresses that gap empirically.
It does not yet claim to measure active secrets at ecosystem scale; the validation campaign remains unfinished and is constrained by a responsible-disclosure protocol.
Instead, the chapter builds the measurement path needed to reach that result.
It first recovers and scans archived contents associated with secret-looking history alterations, establishing how many downloaded blobs contain machine-detectable secret patterns.
It then describes how validity should be tested without invasive access: offline syntactic checks where possible, provider-assisted validation as the preferred path, and researcher-run read-only checks only with explicit platform permission.
The direction of the research is therefore practical as much as descriptive.
If secrets removed from history are still active often enough to matter, the evidence can support better platform protocols: faster provider notification, clearer distinction between cleanup and rotation, and remediation workflows that close the credential rather than only hiding the repository evidence.

\paragraph{What this chapter establishes.}
The chapter first describes the implemented detection and provenance pipeline, from the external \texttt{secrets.pkl} input to Software Heritage content resolution, blob download, gitleaks and TruffleHog scanning, and provenance expansion (\Cref{sec:sec-methodology}).
It then reports the current early detection results: the attrition from filename-filtered candidates to downloaded blobs, and the number of blobs with at least one scanner finding (\Cref{sec:sec-bearing-files}).
Finally, it lays out the non-intrusive validation protocol that would turn those findings into an active-credential study, before stating the threats to validity and the remaining work needed to complete the measurement (\Cref{sec:sec-validation,sec:sec-validity,sec:sec-conclusion}).

\section{Study Design}
\label{sec:sec-methodology}

The study begins from an external prior-work artifact, \texttt{secrets.pkl}, which contains \num{365591843} rows.
Each row records a file path touched in a history-alteration context, with fields for the origin, revision SWHID, branch, before-alteration snapshot, path, status, and basename.
The exact job that produced this dataframe is not versioned in the secret-removal repository; the current work treats it as an input and documents its schema and size.

The implemented pipeline then proceeds in six stages, shown in \Cref{fig:sec-pipeline}.
First, PostgreSQL path filters over the \texttt{altered\_histories.modified\_files} table export sensitive filename candidates.
Second, \texttt{secrets.ipynb} applies a regular-expression filename filter, retaining high-recall patterns such as SSH key names, PEM and key containers, environment files, credential files, token files, secret files, Terraform variable files, and configuration files.
Third, the Rust binary \texttt{extract\_content\_swhids} resolves each retained \texttt{(revision, path)} pair to a Software Heritage content SWHID using the 2025-05-18 graph snapshot.
Fourth, \texttt{swhids\_to\_sha1} maps content SWHIDs to SHA1 digests through the matching digestmap, and \texttt{download\_contents.sh} downloads the raw blobs from the Software Heritage S3 content store.
Fifth, \texttt{content\_lookup} builds an index from content SHA1s back to all matching \texttt{(revision, path)} occurrences.
Finally, \texttt{count\_secrets} scans downloaded blobs with gitleaks and TruffleHog and joins each finding to provenance.

\begin{figure}[t]
\centering
\fbox{%
\begin{minipage}{0.92\linewidth}
\small
\textbf{Figure placeholder (secret-removal detection pipeline).}
Draw the implemented six-stage pipeline: \texttt{altered\_histories} database and external \texttt{secrets.pkl} $\rightarrow$ PostgreSQL path retrieval $\rightarrow$ notebook filename refinement $\rightarrow$ content SWHID resolution with the 2025-05-18 SWH graph $\rightarrow$ SHA1 mapping and S3 blob download $\rightarrow$ SHA1-to-provenance index $\rightarrow$ gitleaks and TruffleHog scan with provenance join.
\end{minipage}%
}
\mycaption[Secret-removal detection pipeline]{%
Placeholder for the implemented secret-removal detection pipeline.%
}
\label{fig:sec-pipeline}
\end{figure}

The units of analysis are nested.
A \emph{candidate pair} is a \texttt{(revision SWHID, path)} tuple retained by the filename filter.
A \emph{content blob} is the deduplicated archived file identified by a content SWHID and SHA1.
A \emph{finding} is a scanner match, represented by the detector tool, secret type, SHA1, and line range, and then expanded to all candidate pairs that reference the same SHA1.
This expansion is necessary to recover provenance, but it also means that expanded finding counts are not the same as unique leaked files.

\section{Early Detection Results}
\label{sec:sec-bearing-files}

The current run is best read as an attrition funnel, not as a completed active-secret study.
From \num{365591843} rows in \texttt{secrets.pkl}, filename refinement retained \num{2846017} candidate \texttt{(revision, path)} pairs.
The provenance index resolved \num{2187481} pairs to content SHA1s, leaving \num{658536} unresolved pairs.
The pipeline identified \num{281740} unique content SWHIDs and SHA1s, of which \num{281728} blobs were downloaded; \num{12} downloads failed.
Among the downloaded blobs, \num{25859} had at least one gitleaks or TruffleHog finding, for an overall detection rate of \SI{9.18}{\percent}.
\NOTEside{interpretation}
This does not mean that \SI{9.18}{\percent} of all altered files in the prior study contained secrets; it is a rate over downloaded blobs after a filename filter that intentionally enriched the sample for likely secret-bearing paths.

\begin{figure}[t]
\centering
\fbox{%
\begin{minipage}{0.92\linewidth}
\small
\textbf{Figure placeholder (current-run attrition funnel).}
Draw a funnel with these stages and counts: \texttt{secrets.pkl} rows \num{365591843}; filename-filtered candidate pairs \num{2846017}; resolved index rows \num{2187481}; unique content SHA1s \num{281740}; downloaded blobs \num{281728}; blobs with at least one finding \num{25859} (\SI{9.18}{\percent} of downloaded blobs).
\end{minipage}%
}
\mycaption[Current-run attrition funnel]{%
Placeholder for the attrition funnel of the current secret-removal pipeline run.%
}
\label{fig:sec-attrition}
\end{figure}

\begin{table}[t]
\centering
\caption{Current pipeline counts for secret detection in history-altered files.}
\label{tab:sec-current-run}
\small
\begin{tabular}{lr}
\hline
\textbf{Stage or artifact} & \textbf{Count} \\
\hline
\texttt{secrets.pkl} rows & \num{365591843} \\
Filename-filtered \texttt{(revision, path)} pairs & \num{2846017} \\
Resolved index rows & \num{2187481} \\
Unresolved candidate pairs & \num{658536} \\
Unique content SWHIDs / SHA1s & \num{281740} \\
Downloaded blobs & \num{281728} \\
Blobs with at least one finding & \num{25859} \\
Overall blob-level detection rate & \SI{9.18}{\percent} \\
gitleaks expanded findings & \num{216850} \\
TruffleHog expanded findings & \num{169862} \\
\texttt{(tool, secret\_type)} summary groups & \num{142} \\
\hline
\end{tabular}
\end{table}

The expanded scanner outputs are larger than the positive-blob count because each finding is joined back to every candidate \texttt{(revision, path)} pair that references the same SHA1.
The current run contains \num{216850} expanded gitleaks rows and \num{169862} expanded TruffleHog rows, and the combined summary contains \num{142} distinct \texttt{(tool, secret\_type)} groups.
The methodology document gives one worked example: the gitleaks \texttt{private-key} row has \num{12009} raw matches spanning \num{11613} unique SHA1s, which expand to \num{145805} distinct \texttt{(revision, path)} locations.
This example illustrates why the chapter reports blob-level, finding-level, and provenance-expanded metrics separately.

These numbers answer only the detection part of the study.
They show that the archived contents of history-altered, credential-looking paths include scanner-detectable secrets.
They do not establish that any detected credential is valid, authorised for a particular service, or still active.
That distinction is central to the rest of the chapter, because the active-rate result is exactly the work that remains unfinished.

\section{Validation Protocol}
\label{sec:sec-validation}

The protocol separates detection from validation.
Detection reports that a blob contains a pattern recognised by gitleaks or TruffleHog.
Validation asks whether the credential can still authenticate to a provider or service.
The protocol defines three tiers so that the study can measure activity without treating third-party credentials as tools for unauthorised access.

\begin{itemize}
\item \textbf{Tier A: offline validation.} Local structural or cryptographic checks issue no network traffic. Examples include format validation, \texttt{ssh-keygen -y}, and OpenSSL parsing. Tier A can reject placeholders, truncated material, and misnamed files, but it cannot prove activity because a revoked credential can remain syntactically valid.
\item \textbf{Tier B: provider-assisted validation.} The researcher contacts the affected provider through an approved security channel and asks the provider to check the credentials on its side, revoke them where appropriate, and return only aggregate activity counts.
\item \textbf{Tier C: researcher-run network validation.} The researcher performs a single read-only or metadata-only check only after explicit written provider authorisation for the precise command, or under a published provider programme that clearly authorises the same class of check.
\end{itemize}

The protocol assigns one predefined check per credential family.
Each check is read-only or metadata-only, and the ``active'' outcome is fixed in advance.
The examples below illustrate the pattern; they are run only after Tier~A filtering and, for network checks, after provider authorisation (Tier~B preferred, Tier~C as fallback).

\begin{itemize}
\item \textbf{SSH private key (GitHub).}
Offline: \texttt{ssh-keygen -y -f ./key} confirms syntactic validity (exit~0).
With authorisation, a single SSH banner probe:
\begin{verbatim}
ssh -i ./key -o IdentitiesOnly=yes -T git@github.com
\end{verbatim}
Active iff the banner contains \texttt{Hi <user>!}.

\item \textbf{GitHub personal access token.}
\begin{verbatim}
curl -H "Authorization: Bearer ghp_..." https://api.github.com/user
\end{verbatim}
Active iff HTTP~200.

\item \textbf{AWS access key.}
\begin{verbatim}
aws sts get-caller-identity
\end{verbatim}
Active iff HTTP~200 with an account ARN returned.

\item \textbf{Slack bot token.}
\begin{verbatim}
curl -X POST -d "token=xoxb-..." https://slack.com/api/auth.test
\end{verbatim}
Active iff the JSON response contains \texttt{"ok":true}.

\item \textbf{Stripe secret key.}
\begin{verbatim}
curl https://api.stripe.com/v1/charges -u sk_live_...:
\end{verbatim}
Active iff HTTP~200 (a metadata-only list request, not a charge creation).

\item \textbf{OpenAI API key.}
\begin{verbatim}
curl -H "Authorization: Bearer sk-..." https://api.openai.com/v1/models
\end{verbatim}
Active iff HTTP~200.
\end{itemize}

These checks do not modify provider state: they authenticate identity or list metadata.
The preferred path remains Tier~B, where the provider performs equivalent checks internally and returns only an aggregate count of credentials that appeared active.

\begin{figure}[t]
\centering
\fbox{%
\begin{minipage}{0.92\linewidth}
\small
\textbf{Figure placeholder (validation decision flow).}
Draw the validation decision flow: scanner finding $\rightarrow$ Tier~A syntactic check $\rightarrow$ provider contact $\rightarrow$ preferred Tier~B provider-assisted aggregate validation or authorised Tier~C read-only check $\rightarrow$ report active/revoked only for authorised families; otherwise report syntactically valid or plausibly exposed with active status unknown.
\end{minipage}%
}
\mycaption[Validation decision flow]{%
Placeholder for the planned non-intrusive validation and responsible-disclosure workflow.%
}
\label{fig:sec-validation-flow}
\end{figure}

The denominator for any future active-rate claim is therefore not the raw detector output.
It is the set of syntactically valid secrets that pass Tier A in a credential family for which Tier B or Tier C validation is authorised.
The numerator is the count of credentials confirmed active by the provider or by the authorised check.
When neither Tier B nor Tier C is available, active status is reported as unknown rather than inferred from silence.

The protocol also excludes several families from active validation.
Generic certificate containers such as \texttt{*.pem}, \texttt{*.key}, \texttt{*.pfx}, and \texttt{*.p12} do not identify a target service by filename alone; if no subject alternative name points to a reachable host, they can only be reported as plausibly exposed.
Database credentials are excluded because probing private hosts through TCP authentication is more intrusive than a metadata-only provider request.
SSH keys without a Git-hosting target are also excluded, because probing arbitrary hosts is not ethically justified by the protocol.

Provider outreach is specified before any Tier C action.
The template asks a provider for a secure channel, states the credential family and count of syntactically valid credentials, and explicitly says that no live validation has yet occurred.
The preferred path is provider-assisted validation and revocation.
If the provider does not respond, the protocol schedules one follow-up at day~14, closes the request at day~30, and sends only a minimal disclosure notice.
Lack of response is not treated as authorisation.

\section{Threats to Validity}
\label{sec:sec-validity}

The first limitation is the external origin of \texttt{secrets.pkl}.
The secret-removal repository documents the dataframe schema and row count, but it does not version the job that created it from the altered-histories database.
The current chapter can therefore describe the downstream pipeline and its current results, but it cannot independently audit the construction of the \num{365591843}-row input within this repository.

The second limitation is selection and coverage.
The filename filter is intentionally high-recall, so many retained paths are configuration or credential-looking files without detector-visible secrets.
At the same time, it can miss secrets stored under unusual names.
Resolution and download also remove data: only \num{2187481} of \num{2846017} candidate pairs resolve into the content lookup index, and \num{12} of \num{281740} unique blobs failed to download.
The detection result is consequently conditioned on successfully resolved and downloaded blobs.

The third limitation is scanner accuracy.
gitleaks and TruffleHog use different detector rules and type taxonomies, so their \texttt{secret\_type} labels are not directly comparable.
They can miss secrets outside their modeled formats, encoded or split credentials, and credentials that require repository context.
They can also flag non-secret high-entropy values or examples.
The planned Tier A validation step is designed to reduce this noise before any active-rate denominator is formed.

The final limitation is the most important one for this chapter: no credential validation has been executed.
\TODO{T-CH5-RESULTS: run or document the Tier A/B/C validation campaign before claiming active, revoked, or rotated credential counts.}
The current results are detection results over archived content.
The study protocol explains how to measure active credentials responsibly, but the repository does not contain outreach logs, provider responses, authorised Tier C runs, or aggregate active-rate outputs.
Any statement about still-active secrets must therefore remain prospective until that campaign exists.

\section{Conclusion and Perspectives}
\label{sec:sec-conclusion}

This chapter reframed the secret-suppression case from Chapter~\ref{chap:histories} as a distinct, unfinished security study.
The implemented contribution is a pipeline that starts from history-altered, credential-looking paths, resolves their archived content in Software Heritage, scans the downloaded blobs with two secret detectors, and links findings back to revision/path provenance.
On the current run, \num{25859} of \num{281728} downloaded blobs have at least one detector finding, a blob-level rate of \SI{9.18}{\percent}.
That result establishes that the archived contents of altered histories contain scanner-detectable secrets.
It does not establish whether those credentials remain active.

The next step is to turn this detection pipeline into an active-secret measurement without crossing ethical or legal boundaries.
The first practical task is to stabilise the current run metadata: record the run date, pin the gitleaks and TruffleHog versions, fill the PostgreSQL export count, and resolve the small discrepancy between \num{281728} downloaded blobs in the results snapshot and the \num{281729} scanned files reported in notebook output.
\TODO{T-CH5-REPRO: record run date, scanner versions, PostgreSQL export count, and resolve the downloaded-vs-scanned blob-count discrepancy before final submission.}
Without those details, the detection result is useful as early evidence but not yet a fully reproducible final experiment.

The second task is to compute Tier A denominators per credential family.
This means applying offline structural checks before any provider contact: parsing private keys, certificates, service-account JSON files, and token formats where such checks are possible.
The output of this step should not be an active-rate result.
It should be a family-by-family count of syntactically valid credentials, placeholders rejected, malformed findings, and families excluded from validation.
That denominator design matters because it prevents raw scanner output from being mistaken for confirmed credentials.

The third task is provider coordination.
The protocol already defines the preferred process: contact providers through a security channel, avoid sending raw credentials in the initial email, request provider-side validation and revocation where possible, and use researcher-run network checks only with written authorisation.
The outreach logs must record dates, channels, follow-up status, closure status, and whether the provider authorised Tier B, authorised Tier C, declined, or did not respond.
The future active-rate result should be reported only for authorised families, together with the number of syntactically valid credentials whose status remains unknown because no authorisation was obtained.

The fourth task is to answer the temporal and comparative questions left open by the protocol.
RQ3 requires pairing the first exposure of a detected secret with the remediation attempt that removed it from the public history.
RQ4 requires comparing rotation or activity outcomes across secret families, but only after the validation campaign supplies active, revoked, and unknown outcomes.
Both analyses should keep the distinction between blobs, findings, and expanded revision/path occurrences, because each denominator answers a different question.

The broader perspective is that secret removal exposes the limit of repository-level remediation.
\NOTEside{interpretation}
The current evidence shows that history rewriting can remove a credential from the visible repository state while leaving the archived content available for later analysis.
The unfinished part is to measure how often this archival exposure remains operationally dangerous because the credential itself was not revoked.
Completing that measurement would turn the chapter from an early detection study into a direct evaluation of whether developers and providers close the loop after secret exposure.
